\documentclass[11pt,a4paper]{article}

\usepackage[T1]{fontenc}
\usepackage[utf8]{inputenc}
\usepackage[margin=2.2cm]{geometry}
\usepackage{xcolor}
\usepackage{longtable}
\usepackage{array}
% enumitem is not part of BasicTeX; these keep lists compact without it, so the
% document rebuilds on a stock TeX installation with no extra packages.
\newcommand{\tightlist}{\setlength{\itemsep}{0pt}\setlength{\parskip}{0pt}\setlength{\parsep}{0pt}}
\usepackage{fancyhdr}
\usepackage[hidelinks]{hyperref}

\definecolor{sev-high}{rgb}{0.70,0.10,0.10}
\definecolor{sev-med}{rgb}{0.80,0.45,0.00}
\definecolor{sev-low}{rgb}{0.25,0.45,0.25}
\definecolor{codebg}{rgb}{0.96,0.96,0.94}

\newcommand{\sevhigh}{\textcolor{sev-high}{\textbf{HIGH}}}
\newcommand{\sevmed}{\textcolor{sev-med}{\textbf{MEDIUM}}}
\newcommand{\sevlow}{\textcolor{sev-low}{\textbf{LOW}}}
\newcommand{\confirmed}{\textcolor{sev-low}{\textbf{CONFIRMED}}}
\newcommand{\unverified}{\textcolor{sev-med}{\textbf{UNVERIFIED}}}

\newcommand{\finding}[2]{\subsection*{#1 \hfill \normalsize #2}\addcontentsline{toc}{subsection}{#1}}

\pagestyle{fancy}
\fancyhf{}
\lhead{MOSIP Decode --- PS1 Findings Log}
\rhead{\thepage}
\renewcommand{\headrulewidth}{0.4pt}

\title{\textbf{MOSIP Decode 2026 --- Problem Statement 1}\\[0.3em]
\large Automated Conformance Testing for Inji Certify \& Inji Verify\\[0.8em]
\normalsize Findings Log and Open Questions}
\author{Shardul Chogale \and Mukta Varak}
\date{Document started: 20 September 2026\\Last updated: 2 October 2026}

% Long code identifiers (\texttt) otherwise never break and run into the margin. Allow
% breaks at hyphens inside them, and a little extra inter-word stretch as a last resort.
\let\origtexttt\texttt
\renewcommand{\texttt}[1]{\origtexttt{\hyphenchar\font=`\-\relax #1}}
\setlength{\emergencystretch}{3em}
\renewcommand{\_}{\textunderscore\allowbreak}
\begin{document}
\maketitle
\thispagestyle{fancy}

\begin{abstract}
\noindent This document records, day by day, the conformance testing work carried out
against the Inji stack as part of MOSIP Decode Problem Statement 1. It captures what was
run, what failed, the evidence for each failure, and the questions we need answered by
MOSIP mentors. Every finding is tagged as \confirmed{} or \unverified{}; we do not report
a defect upstream until it is reproducible and the root cause is understood.
\end{abstract}

\tableofcontents
\newpage

%=========================================================
\section{Scope and System Under Test}

\subsection*{Problem statement}
Build an automated harness that spins up Inji Certify and Inji Verify, drives the OpenID
Foundation conformance suite's REST API programmatically to run the issuer and verifier
test plans, and feeds the results into each module's existing \texttt{api-testrig},
producing benchmark-gated reports with CI/CD integration.

\subsection*{Components under test}
\begin{longtable}{p{4.5cm} p{4cm} p{6cm}}
\hline
\textbf{Component} & \textbf{Version} & \textbf{Notes}\\
\hline
\endhead
Inji Verify (service) & 0.18.2 & \texttt{injistack/inji-verify-service}\\
Inji Verify (UI) & 0.18.2 & \texttt{injistack/inji-verify-ui}, amd64 under emulation on arm64\\
Inji Certify & 0.14.0 & Not yet exercised (deferred, see Plan)\\
OpenID Conformance Suite & 5.3.1 (rev 440eec8) & Local Docker deployment, prebuilt images\\
\hline
\end{longtable}

\subsection*{Test environment}
\begin{itemize}\tightlist
  \item Host: Apple Silicon (arm64) macOS, Docker Desktop.
  \item Conformance suite reachable at \texttt{https://localhost.emobix.co.uk:8443}
        (self-signed certificate; hostname resolves publicly to 127.0.0.1).
  \item Inji Verify exposed publicly via an ngrok HTTPS tunnel. A public hostname is
        mandatory because Verify identifies itself with a \texttt{did:web} DID, which is
        resolved over public HTTPS.
  \item The conformance suite's REST API required no authentication token in the local
        \texttt{dev} profile, consistent with the problem statement.
\end{itemize}

\subsection*{Evidence handling}
Conformance suite logs are exported as signed JSON (\texttt{.json} plus \texttt{.json.sig})
and stored under \texttt{logs/<date>/}. Findings below cite the exact log line
(\texttt{src} field) and the OID4VP specification section reported by the suite.

\newpage
%=========================================================
\section{Day Log}

\subsection{Day 1 --- Saturday, 20 September 2026}

\subsubsection*{Objective}
Stand up the conformance suite and one MOSIP component, and complete one manual
conformance test end to end to establish a baseline.

\subsubsection*{Decisions taken}
\begin{enumerate}\tightlist
  \item \textbf{Inji Verify before Inji Certify.} Verify's Docker Compose brings up three
        containers with no external dependencies. Certify's \texttt{docker-compose-injistack}
        requires a PKCS12 keystore obtained through an onboarding process, a public DID
        endpoint, an external authorisation server, and partner API keys. The verifier test
        plan also has 12 modules against the issuer plan's 61. Verify is the cheaper path to
        a first result, and everything learned transfers.
  \item \textbf{Non-HAIP plan first.} The HAIP verifier plan offers only
        \texttt{direct\_post.jwt} as a response mode. Inji Verify hardcodes
        \texttt{direct\_post} (see F-06), so no HAIP test can complete. We ran
        \emph{OpenID for Verifiable Presentations 1.0 Final: Test a verifier --- alpha tests}
        instead, which offers plain \texttt{direct\_post}.
  \item \textbf{Credential format \texttt{sd\_jwt\_vc}.} The plan offers only
        \texttt{sd\_jwt\_vc} and \texttt{iso\_mdl}. Verify has no mdoc support whatsoever
        (F-08), making \texttt{iso\_mdl} unusable.
\end{enumerate}

\subsubsection*{Test plan configuration}
\begin{itemize}\tightlist
  \item Plan: \texttt{oid4vp-1final-verifier-test-plan}, plan ID \texttt{lN7C4DH1HvYmc}
  \item Variant: \texttt{credential\_format=sd\_jwt\_vc}, \texttt{client\_id\_prefix=redirect\_uri},
        \texttt{request\_method=url\_query}, \texttt{vp\_profile=plain\_vp},
        \texttt{response\_mode=direct\_post}
  \item Alias: \texttt{injiverify-shardul}
\end{itemize}

\subsubsection*{Runs executed}
\begin{longtable}{p{2.2cm} p{3.4cm} p{8.5cm}}
\hline
\textbf{Test ID} & \textbf{Outcome} & \textbf{Notes}\\
\hline
\endhead
\texttt{dgeeosZpIfDK1bA} & FAILED / INTERRUPTED & First attempt. Request object had expired
(5-minute window). Discarded as an artefact of manual timing, not a defect.\\
\texttt{dMg4B5whxWgShIB} & FAILED / INTERRUPTED & Fresh request (2m05s old). Suite never
dereferenced \texttt{request\_uri} because the \texttt{url\_query} variant expects inline
parameters. Root cause understood; see F-09.\\
\texttt{MtzpWoD2dVgtic0} & FAILED / INTERRUPTED & After switching the credential's
\texttt{clientIdScheme} to \texttt{pre\_registered}. Reached 20+ checks; produced the
substantive findings below. \textbf{This is the Day 1 baseline run.}\\
\hline
\end{longtable}

\subsubsection*{Configuration change made}
In \texttt{inji-verify/docker-compose/config/config.json}, the credential
\emph{Mock Identity (SD JWT)} was changed from \texttt{"clientIdScheme":"did"} to
\texttt{"clientIdScheme":"pre\_registered"}. Per Inji Verify's own README, \texttt{did}
delivers the authorisation request \emph{by reference} via \texttt{request\_uri}, whereas
\texttt{pre\_registered} delivers it \emph{inline}. Only the inline form matches the
\texttt{request\_method=url\_query} variant. Backup retained at \texttt{config.json.bak}.

\subsubsection*{Result of the baseline run (\texttt{MtzpWoD2dVgtic0})}
Approximately fourteen checks passed, including nonce extraction, response type,
response mode, absence of \texttt{client\_id\_scheme}, absence of \texttt{scope}, absence of
\texttt{redirect\_uri}, response URI validity, and the outer parameter allow-list. Three
checks failed and three raised warnings. The test terminated at
\texttt{ExtractDCQLQueryFromAuthorizationRequest}.

\subsubsection*{Status at end of Day 1}
Test 01 (\texttt{happy-flow}) complete and documented. Tests 02--11 deferred.

\subsection{Day 2 --- Tuesday, 23 September 2026}

\subsubsection*{Calendar note}
The work recorded above took place on 20 September, \emph{three days before MOSIP Decode
officially began}. The event kicked off on 23 September; the planned kickoff seminar and
mentor AMA session did not take place. No work was carried out between 20 and 23 September,
so everything listed as ``deferred to Day 2'' above remained outstanding at the start of
the event proper. This is recorded because the project plan was originally dated from
20 September and therefore measured against the wrong calendar; it has since been
re-anchored in \texttt{docs/PLAN.md}.

The practical consequence is favourable: the environment and a first evidenced baseline
existed on day two of a four-week event, and the questions in Section 4 were already
formulated before any mentor contact was possible.

\subsubsection*{Objective}
Planning and documentation only. No conformance runs were executed, so no findings are
added or amended in this entry.

\subsubsection*{Decisions taken}
\begin{enumerate}\tightlist
  \item \textbf{Deliverable code location.} The conformance runner, plan configurations and
        orchestration live in our own repository. The \texttt{OpenIDConformanceTest} classes
        and their wiring go on branches of our forks of \texttt{inji-verify} and
        \texttt{inji-certify}, shaped so that each could be opened as an upstream pull
        request. Note that \texttt{api-test/} lives inside those two repositories, not in
        \texttt{mosip-functional-tests}, which carries only \texttt{apitest-commons}.
  \item \textbf{Target specification treated as configuration.} Because Q1 has no channel
        for an answer other than the community forum, and may not receive one, the runner
        will take the plan name and variant set as configuration rather than hard-coding
        them. Either answer to Q1 then costs a configuration change rather than a rewrite.
  \item \textbf{Benchmark semantics fixed provisionally as ``regression from recorded
        baseline''} (option (a) in Q4), with an expected-failures file under version
        control. Given that the current honest baseline is a full set of failures at DCQL
        extraction, an absolute pass-rate gate would be meaningless.
\end{enumerate}

\subsubsection*{Status at end of Day 2}
No change to the findings catalogue. Modules 02--11, the ID2 verifier plan, and the
resolution of F-05 remain outstanding and are scheduled within Week 1.

\subsection{Day 3 --- Wednesday, 24 September 2026}

\subsubsection*{Objective}
Build the Week 2 conformance runner, and before writing any of it, settle the two
unknowns that \texttt{docs/WEEK2.md} identified as determining the week's shape. Both were
resolved by reading the conformance suite's source in the local clone rather than by
experiment.

\subsubsection*{Unknown 1 resolved --- the verifier handoff needs no browser}

The verifier plans are inverted relative to the issuer plans: the suite plays the wallet and
waits passively for the verifier to initiate. Its web UI presents this as a box into which a
tester pastes an \texttt{openid4vp://} URI, which is why every run so far has been manual.
The open question was whether that paste has an API equivalent, because if it did not, the
entire week would have had to be re-planned around a semi-automated flow.

It does. \texttt{AbstractVP1FinalVerifierTest.start()} reads:

\begin{quote}\small
\texttt{getBrowser().requestUriInput(env.getString("authorization\_endpoint"),}\\
\texttt{\ \ "Paste the openid4vp:// authorization request produced by the verifier under}\\
\texttt{\ \ test; its query string will be delivered to this test's authorization endpoint.");}
\end{quote}

\noindent and \texttt{handleHttp()} dispatches on \texttt{path.equals("authorize")}. The test
exposes \texttt{authorization\_endpoint} via \texttt{exposeEnvString}, readable through
\texttt{GET /api/runner/\{id\}}. Delivering an authorisation request is therefore an ordinary
HTTP GET to that endpoint carrying the verifier's query parameters. The paste box is a
convenience wrapper over exactly that call.

\textbf{Consequence.} A fully unattended verifier run is achievable. This substantially
answers Q6, which we had expected to need a mentor.

\subsubsection*{Unknown 2 resolved --- the screenshot gate is configuration, not manual work}

The happy-flow module requires a screenshot of the verifier's verification-result page before
reaching REVIEW. \texttt{AbstractVP1FinalVerifierTest} carries a purpose-built path for
automated runs: \texttt{handleVerificationEvidenceRequest()} serves an HTML stand-in page,
and \texttt{fillScreenshotPlaceholderViaBrowserAutomationIfConfigured()} fills the placeholder
when the test configuration contains a \texttt{browser} automation entry matching the
evidence URL. The source comment is explicit that this exists because ``in automated runs
there is no verifier UI a human could take a real screenshot of''.

A \texttt{browser} block has been added to \texttt{configs/runner.json} accordingly.
\textbf{This is unconfirmed against a live run} --- the match pattern in particular is a
reasonable guess at the evidence URL and must be validated before it is relied upon.

\subsubsection*{Work completed}
\begin{itemize}\tightlist
  \item \texttt{runner/} package: configuration loading, a conformance suite API client, an
        Inji Verify client, the handoff, result normalisation, the benchmark gate, an
        orchestrator and a CLI entry point.
  \item \texttt{run-conformance.sh} providing \texttt{--component verify | --component
        certify | --combined}, as the problem statement requires.
  \item 24 unit tests, all passing, covering normalisation against the real Day 1 log and
        the gate's behaviour --- including its failure path, since a gate that cannot fail
        is decoration.
\end{itemize}

\subsubsection*{Defect found in our own code}
The gate initially set \texttt{regression} and \texttt{improvement} only on the branches
that made them true, relying on the normalisation layer having pre-populated the defaults.
Applied to any module assembled by another path this raised \texttt{KeyError} rather than
reporting ``no regression''. Caught by \texttt{test\_gate.py} before any live run and fixed
by setting both unconditionally. Recorded because it is the argument for writing the gate's
failure-path tests first: the bug lived in the success path and only the absent-key case
exposed it.

\subsubsection*{Status at end of Day 3}
Runner complete and unit-tested, but \textbf{not yet executed against a live conformance
suite} --- Docker was not running on the development machine during this session. Nothing in
this entry should be read as evidence that the harness works end to end. The first live run
is the next task, and until it happens the handoff and screenshot mechanisms are
source-derived expectations rather than confirmed behaviour.

No new findings against Inji Verify; no conformance modules were executed.

\subsection{Day 4 --- Monday, 29 September 2026}

\subsubsection*{Objective}
First execution of the harness against a live conformance suite. The Day 3 runner had been
unit-tested only, and every claim about the handoff was source-derived. Today's question was
whether the automated run measures the same thing as the manual Day 1 run --- not merely
whether it completes.

\subsubsection*{Environment}
Conformance suite 5.3.1 (prebuilt images), Inji Verify 0.18.2, ngrok 3.39.11 on the static
domain. Both uncommitted local modifications (the ngrok host in \texttt{docker-compose.yml},
\texttt{pre\_registered} for \emph{Mock Identity (SD JWT)}) confirmed intact after five days.
Same variant as Day 1.

\subsubsection*{Runs executed}
\begin{longtable}{p{2.6cm} p{2.4cm} p{9.2cm}}
\hline
\textbf{Test ID} & \textbf{Outcome} & \textbf{Notes}\\
\hline
\endhead
\emph{(none)} & HARNESS ERROR & Run 1. Suite side worked; Inji Verify returned HTTP 400 to
\texttt{/vp-session-request}. Our defect, not Verify's --- see below. No module result.\\
\texttt{hcuWEvxGbvRgt7e} & FAILED / INTERRUPTED & Run 2, \texttt{nonceMode=sdk}. \textbf{First
fully unattended run.} Reproduces the Day 1 manual baseline exactly.\\
\texttt{Q9MaDZyfpJhjiDp} & FAILED / INTERRUPTED & Run 3, \texttt{nonceMode=service}. Controlled
experiment; differs from run 2 in the nonce source only.\\
\hline
\end{longtable}

\subsubsection*{Result 1 --- the automated run reproduces the manual baseline exactly}
Run 2 and the Day 1 manual baseline (\texttt{MtzpWoD2dVgtic0}) produced identical counts
--- 15 success, 4 failure, 3 warning, 7 info --- and the identical set of seven non-success
checks, terminating at the same place, \texttt{ExtractDCQLQueryFromAuthorizationRequest}. No
check appears in one run and not the other. The handoff delivered with the full 300 seconds
remaining, and the suite answered its \texttt{/authorize} endpoint with HTTP 302.

This is the result that validates the harness. Completing is not enough; a harness that
completed while measuring something different from what a tester would observe by hand
would be worse than none. The Day 3 expectation that delivery is an ordinary HTTP GET to
the exposed \texttt{authorization\_endpoint} is now confirmed behaviour, not inference.

\subsubsection*{Result 2 --- controlled experiment confirms the F-01/F-02 root cause}
On Day 2 the weak nonce was traced to \texttt{btoa(Date.now())} in the SDK, with
\texttt{verify-service}'s own \texttt{SecureRandom} generator reached only as a fallback. That
was a reading of the source. Runs 2 and 3 test it: they differ in one variable --- whether
the harness supplies the SDK-shaped nonce or lets \texttt{verify-service} generate one.

\begin{longtable}{p{5.4cm} p{4.2cm} p{4.2cm}}
\hline
\textbf{Check} & \textbf{sdk (run 2)} & \textbf{service (run 3)}\\
\hline
\endhead
\texttt{VP1FinalEnsureMinimumNonceEntropy} & WARNING, 71.68 bits & SUCCESS\\
\texttt{CheckForInvalidCharsInNonce} & FAILURE (\texttt{=}) & SUCCESS\\
Nonce sent & \texttt{MTc5MDY5NjY0NTQzNw==} & \texttt{78cb380cda818ebc189b80704e6f17a3}\\
Counts (success / failure / warning) & 15 / 4 / 3 & 17 / 3 / 2\\
\hline
\end{longtable}

Exactly the two nonce checks move; every other check is unchanged and both runs still
terminate at DCQL extraction. F-01 and F-02 are therefore confirmed by three independent
routes: the suite flags them, the source locates them in the SDK, and removing only the
SDK's nonce clears them. The service-generated nonce is 32 hexadecimal characters --- the
16 \texttt{SecureRandom} bytes of \texttt{SecurityUtils.generateNonce()}, as the source
predicted. This materially strengthens the upstream report.

\subsubsection*{Measurement decision recorded: the nonce source}
Driving \texttt{verify-service} directly rather than through the UI changes what is measured.
If the harness sends no nonce, the service generates a sound one and the nonce findings
silently vanish --- not because anything was fixed, but because the harness bypassed the
component that causes them. The runner therefore makes the nonce source explicit
(\texttt{nonceMode}), defaults to \texttt{sdk} so that runs measure what the shipped product
sends a wallet, and records the mode on every module's output. The SDK-shaped nonce is a
deliberate, documented reproduction of the SDK line, verified against the three Day 1
nonces by a unit test.

\subsubsection*{Defects found in our own code}
Four, all fixed today. Recorded because each would have produced wrong evidence.
\begin{enumerate}\tightlist
  \item \textbf{The harness sent no presentation definition.} \texttt{VPRequestCreateDto} takes it
        inline. Found by reading the DTO before the first run.
  \item \textbf{The harness sent the raw config definition, which has no \texttt{id}.}
        \texttt{VPDefinitionResponseDto} requires one, so \texttt{verify-service} answered with
        a bare 400 (run 1). The UI wraps selected credentials in an envelope in
        \texttt{vpVerificationState.ts} (lines 58--63, 82--83): a fixed id, a fixed purpose,
        and \emph{only} the credentials' \texttt{input\_descriptors}. The harness now builds
        the same envelope.
  \item \textbf{The gate passed a run that measured nothing.} Run 1's harness error had no
        result; the severity ordering read that as FAILED, equal to the FAILED baseline, and
        the gate reported PASSED. A harness error now always fails the gate while the module
        itself stays SKIP, so the fault is never recorded against MOSIP. \emph{This is the
        most serious of the four}: with the current baseline, the harness could have broken
        completely and CI would have stayed green.
  \item \textbf{A refactor deleted an exception class}, breaking the CLI while all unit tests
        passed --- none imported the modules that used it. An import test now covers every
        runner module.
\end{enumerate}

\subsubsection*{Observation --- not a finding}
The presentation definition id is a hard-coded constant in
\texttt{vpVerificationState.ts:59}: every request from every deployment of this UI carries
\texttt{c4822b58-7fb4-454e-b827-f8758fe27f9a}. Presentation Exchange describes the id as
identifying the definition, and a constant is arguably fine for an unchanging definition ---
but here the definition's contents vary with which credentials are selected while the id
does not. Not catalogued: we have not established that anything relies on the id being
unique, and without a consequence it is not a defect. Raised as a question instead.

\subsubsection*{Status at end of Day 4}
The harness runs unattended end to end and reproduces the manual baseline check-for-check.
Carried-forward Week 1 items remaining: modules 02--11 and the ID2 verifier plan (now a
single command each), and F-05, which today's runs cannot settle --- the harness forwards
\texttt{client\_metadata} only when Verify supplies it, and in the \texttt{pre\_registered}
path it does not, so the open question of whether the UI adds it at QR-build time still
needs the UI's QR payload. The screenshot mechanism is not yet exercised: both runs
terminated at DCQL before reaching it.

Automated logs are archived unsigned under \texttt{logs/2026-09-29/}; anything reported
upstream still needs the suite's signed export.

\subsection{Day 5 --- Wednesday, 30 September 2026}

\subsubsection*{Objective}
Prepare for the first mentor contact, and turn the project's documentation into a single
starting point.

\subsubsection*{Facts established}
\begin{itemize}\tightlist
  \item \textbf{The final submission deadline is 26 October 2026, 11:59 PM IST}, read from
        the official event page. Every plan so far assumed about 21 October; the real date
        adds five days, now used as a fifth week for finishing and submitting.
  \item \textbf{Mentor contact is the Weekly Connect}: Wednesdays, 5:00--6:00 PM IST, from
        today (from the organisers' email of 22 September). Planned checkpoints follow it:
        30 September, 7, 14 and 21 October.
  \item Submission requires an overview plus a 1--2 page document or a 5--7 slide pitch;
        multiple submissions are allowed and the last before the deadline counts. The
        detailed submission and evaluation criteria sit behind separate links and are to
        be confirmed at the session.
\end{itemize}

\subsubsection*{Work completed}
\texttt{docs/handbook/handbook.tex}: a mentor brief ranking the open questions by how much
each answer changes the build, with our default for each; a start-to-end guide from a
clean machine to shutdown and submission; a record of every step taken since 16
September; and a day-by-day planner for both of us to 26 October. \texttt{PLAN.md},
\texttt{CONTEXT.md} and \texttt{WEEK2.md} updated to the real deadline.

\subsubsection*{New question raised}
The problem statement asks for Inji Certify and Inji Verify ``on 1.0'', with ``1.0.x''
repositories. The newest images we could run are Verify 0.18.2 and Certify 0.14.0. If ``1.0''
names a release line we are not running, every result so far is against the wrong
version. Added to the mentor brief as the second-ranked question.

\subsubsection*{Status at end of Day 5}
No conformance runs today; no findings added or changed. No reply yet from Mukta to the
29 September setup email. Mentor answers to be recorded here after the session.

\subsection{Day 6 --- Friday, 2 October 2026}

\subsubsection*{Objective}
Absorb the recording of the 23 September kickoff webinar, and act on anything in it that
changes the build.

\subsubsection*{Correction to earlier entries}
Earlier entries, and \texttt{PLAN.md}, said the kickoff seminar did not take place. \textbf{It
did}: the organisers held it on 23 September, 5:00--6:00 PM IST, and recorded it. We did not
attend live, and worked from the recording's transcript today. The ``AMA'' is the weekly
Wednesday session, 5:00--6:00 PM IST, from 30 September; organisers' contacts are
\texttt{decode@mosip.io} and the forum at community.mosip.io.

\subsubsection*{What the kickoff said about PS1}
Walkthrough at 24:58--32:45 of the recording, plus the general session. Points that bear on
our build:
\begin{itemize}\tightlist
  \item \textbf{Mandatory submissions} (08:01): source code, a working prototype, \emph{a slide
        deck, and a video demo}. Stronger than the event page's ``document \emph{or} slides''.
        Settles part of mentor-brief question 6: the video is required.
  \item \textbf{Docker Compose must itself support both modes} (28:30): one file or step that
        brings up each module on its own, or both together, connected to the conformance
        suite. This is a concrete design input for the compose work: selective bring-up,
        for example compose profiles, not one monolithic stack.
  \item \textbf{One-command execution is the core ask} (29:20--29:46): explicitly contrasted with
        configuring the plan, clicking continue and selecting tests in the suite's UI. The
        runner meets this for Verify.
  \item \textbf{CI/CD is framed as good-to-have, with a stated purpose} (27:33--28:26): run the
        suite on every feature change so non-conformance is caught when the change is made,
        not at the end. This is the case for our regression-from-baseline gate.
  \item \textbf{Result diff is motivated by run-to-run variation} (30:12--30:37): ``each run
        might have different failures''. Supports building the diff on module- and
        check-level results, which the contract already carries.
  \item The resources point to ``the latest release 1.0.x'' documentation (32:08). Taken with
        the problem statement's ``on 1.0'', this prompted the version check below.
\end{itemize}

\subsubsection*{Finding --- an Inji 1.0 line exists, and most of our findings do not apply to it}
The Interop Playground walkthrough (41:08) mentions choosing between a ``presentation
exchange request or a DCQL request'' --- implying some Inji Verify sends DCQL, which
contradicts F-03. Checked against upstream rather than taken on trust:

\begin{itemize}\tightlist
  \item \texttt{mosip/inji-verify} has a \texttt{release-1.0.x} branch, under active development
        (head \texttt{6a6ab3b}, 1 October: ``Rename DCQL error codes to snake case''), versions
        \texttt{1.0.0-alpha.2-SNAPSHOT}, and tags \texttt{v1.0.0-alpha.1} (4 August) and
        \texttt{v1.0.0-dev.*}. \texttt{mosip/inji-certify} also has \texttt{release-1.0.x}.
  \item Published on Docker Hub: \texttt{injistack/inji-verify-service} and
        \texttt{-ui} at \texttt{1.0.0-alpha.1}. The branch's compose file names
        \texttt{1.0.0-alpha.2}, which is not yet published.
  \item \textbf{Everything we have run was 0.18.2}, the newest image in the \texttt{0.18.x} line
        and the one the default branch's compose file pinned on 18 September.
\end{itemize}

Our findings against the 1.0 line, from source (not yet from runs):

\begin{longtable}{p{2.4cm} p{3.9cm} p{3.9cm} p{4cm}}
\hline
\textbf{Finding} & \textbf{0.18.2 (tested)} & \textbf{1.0.0-alpha.1 (runnable)} & \textbf{release-1.0.x (head)}\\
\hline
\endhead
F-01/F-02 nonce & \texttt{btoa(Date.now())} & \texttt{crypto.getRandomValues} & \texttt{crypto.getRandomValues}\\
F-03 DCQL & absent & present (16 service files) & present\\
F-06 response mode & \texttt{direct\_post} only & \texttt{direct\_post} only & not determined by source search\\
F-09 client id & \texttt{did}, \texttt{pre\_registered} & \texttt{redirect\_uri:} works --- see runs below$^\dagger$ & \texttt{x509\_san\_dns} configured\\
\hline
\end{longtable}

\textbf{Consequences.}
\begin{enumerate}\tightlist
  \item \textbf{The F-01/F-02 upstream issue is on hold.} As drafted it reports a defect the 1.0
        line has already fixed. It may still be worth raising for the 0.18.x line if that line
        is maintained --- a question for the mentors, not an assumption. This is the
        ``confirm before claiming'' rule doing its job a third time.
  \item \textbf{The baseline should move to \texttt{1.0.0-alpha.1}.} It is the newest runnable 1.0
        build and contains DCQL, so the OID4VP 1.0 Final plan may now produce real passes rather
        than a uniform failure at DCQL extraction. This would also give the gate something to
        protect, which it currently lacks.
  \item \textbf{Every finding in the catalogue is now scoped to 0.18.2}, explicitly, below.
        They remain true and evidenced for that version.
  \item The runner's \texttt{sdk} nonce mode reproduces the 0.18.x SDK. Against 1.0 builds
        the faithful reproduction is a random nonce; the mode needs to follow the version
        under test rather than be a fixed default.
\end{enumerate}

$^\dagger$\,Corrected the same day. Written from source as ``no overlap'' because alpha.1 has no
\texttt{x509\_san\_dns}; run \texttt{5KS13qvmbwOMNtY} below showed the \texttt{redirect\_uri:}
client identifier prefix works, which source search had not revealed. Recorded rather than
silently fixed.

\subsubsection*{Runs on Inji Verify 1.0.0-alpha.1}

Run from a worktree of \texttt{inji-verify} at tag \texttt{v1.0.0-alpha.1}
(\texttt{../inji-verify-1.0}) with the same suite and variant as every earlier run. The runner
gained a 1.0 mode first: \texttt{/v2/vp-session-request}, a DCQL query read from Verify's own
config, and a reproduction of the 1.0 SDK, including the \texttt{client\_metadata} it adds
when building the request URI. Each run changed one thing, so each change in result has a
single cause.

\begin{longtable}{p{2.6cm} p{4.4cm} p{3.4cm} p{4.1cm}}
\hline
\textbf{Test ID} & \textbf{Change from previous} & \textbf{Result} & \textbf{Stopped at / outcome}\\
\hline
\endhead
\texttt{Bob121tXCWTGf9y} & First 1.0 run; shipped \emph{Mock Identity} credential
& INTERRUPTED / FAILED, 23 pass & \texttt{CreateSdJwtKbCredential}: our plan config lacked a
signing key. \textbf{DCQL extracted; all nonce checks pass.}\\
\texttt{gIA512EeJWehcHt} & Test signing key added & FINISHED / FAILED, 28 pass & Verify
answered the submission with \textbf{HTTP 500} --- F-10\\
\texttt{tpKVz6upccxpO5U} & Database init script corrected & FINISHED / FAILED, 28 pass &
Verify answered 400: credential \texttt{vct} does not match the query --- correct behaviour\\
\texttt{ZQqnbYfw46wpmyI} & Credential requesting \texttt{urn:eudi:pid:1} as \texttt{dc+sd-jwt}
& FINISHED / FAILED, 32 pass & \textbf{Verify accepted the presentation (HTTP 200).} Two
DCQL failures cleared. Remaining failures: client id pairing only\\
\texttt{5KS13qvmbwOMNtY} & Client id \texttt{redirect\_uri:<response\_uri>} &
\textbf{FINISHED / REVIEW, 35 pass, 0 fail, 0 warn} & Every automated check passes; screenshot
placeholder filled by the suite's browser automation\\
\hline
\end{longtable}

\textbf{The happy-flow module passes against Inji Verify 1.0.0-alpha.1, fully unattended.}
REVIEW is the suite's state for ``all checks passed, image evidence awaiting human review'',
which this module reaches by design. The suite log confirms ``All placeholders filled by
scripted browser'' --- the screenshot requirement (Q6) is automated, now observed rather than
read from source.

\textbf{What a passing run requires} --- each item established by the run that needed it:
\begin{enumerate}\tightlist
  \item A signing key for the suite's test credential (\texttt{credential.signing\_jwk}),
        generated with the suite's own \texttt{generate-vp-test-cert.py}; kept out of git.
  \item The database fix for F-10.
  \item A verifier credential requesting \texttt{vct} \texttt{urn:eudi:pid:1} as
        \texttt{dc+sd-jwt}. The suite's verifier tests \emph{always} present that credential
        type --- hard-coded in \texttt{CreateSdJwtKbCredential.java:17} --- so a verifier asking
        for anything else will, correctly, reject it. Run \texttt{tpKVz6upccxpO5U} is Inji Verify
        doing exactly that. Added locally to Verify's \texttt{config.json}.
  \item The \texttt{redirect\_uri:} client identifier prefix: \texttt{client\_id} =
        \texttt{redirect\_uri:} followed by the response URI, matching the plan's
        \texttt{client\_id\_prefix=redirect\_uri} variant. Per-developer, because the response
        URI contains the ngrok host.
\end{enumerate}

\textbf{The gate now has something to protect.} The 1.0 baseline records happy-flow as
REVIEW; a return to FAILED on the target version fails the build --- which, while every
module failed, the gate could never do. Baselines are now per version
(\texttt{expected-failures.json} for 1.0; \texttt{expected-failures-0.18.json} for the
earlier evidence).

\subsubsection*{The full plan on Inji Verify 1.0.0-alpha.1}

Plan \texttt{nupFFSGqySsbb}: all 11 modules the OID4VP 1.0 Final verifier plan contains for this
variant (\texttt{invalid-session-transcript} is mdoc-only). Every module FINISHED with 0
failures and 0 warnings: 4 PASSED, 7 REVIEW.

\textbf{REVIEW was not taken at face value.} For the negative tests the suite presents a forged
credential and asks a human to confirm the verifier rejected it. It cannot see this itself:
Inji Verify answers the \texttt{direct\_post} with HTTP 200 and verifies afterwards, which OID4VP
1.0 Final permits. Our screenshot automation fills the placeholder regardless --- so ``REVIEW, 0
failures'' would have been reported even if Verify had \emph{accepted} a forged credential. The
first full run (plan \texttt{YyUQFGfA9rlMr}) had exactly that blind spot.

The harness now settles each REVIEW with evidence only the verifier side has: Inji Verify's
own verdict from \texttt{/vp-session-results}, compared with what the suite's REVIEW message
says the test requires. A contradiction fails the gate. Second full run:

{\footnotesize
\begin{longtable}{>{\raggedright\arraybackslash}p{4.4cm} p{2.95cm} p{1.3cm} >{\raggedright\arraybackslash}p{5.9cm}}
\hline
\textbf{Module} & \textbf{Test ID} & \textbf{Suite} & \textbf{Inji Verify's verdict}\\
\hline
\endhead
happy-flow & \texttt{FiEKKnONeoNFVWu} & REVIEW & accepted; required accept --- consistent\\
minimal-cnf-jwk & \texttt{eGNvPaFZCvXDyF2} & REVIEW & accepted; required accept --- consistent\\
request-uri-method-post & \texttt{D9k8YD8bC5KOPYk} & REVIEW & accepted --- \textbf{not exercised} (n/a to \texttt{url\_query})\\
request-uri-fetched-twice & \texttt{1m9l4wYovOH7Qq5} & REVIEW & accepted --- \textbf{not exercised} (n/a to \texttt{url\_query})\\
invalid-kb-jwt-signature & \texttt{V8O4tr84SWq2nRX} & REVIEW & rejected; required reject --- consistent\\
invalid-credential-signature & \texttt{R7r9vitOs0OMGiG} & REVIEW & rejected; required reject --- consistent\\
invalid-sd-hash & \texttt{jP5ItCI3EM9AT2w} & REVIEW & rejected; required reject --- consistent\\
invalid-kb-jwt-nonce & \texttt{79k2ff0JPhb1Db6} & PASSED & rejected at submission (HTTP 400)\\
invalid-kb-jwt-aud & \texttt{ACAbYUZ9qjIYl6E} & PASSED & rejected at submission (HTTP 400)\\
kb-jwt-iat-in-past & \texttt{wXtLgQRsyy2VjF2} & PASSED & rejected at submission (HTTP 400)\\
kb-jwt-iat-in-future & \texttt{LJnVhzjO72E8sEe} & PASSED & rejected at submission (HTTP 400)\\
\hline
\end{longtable}}

\textbf{Result.} On Inji Verify 1.0.0-alpha.1 the verifier accepted every valid presentation and
rejected every invalid one --- three forged ones (credential signature, key-binding signature,
SD hash) after accepting the submission, and four (nonce, audience, issued-at past and future)
at submission with precise error messages. \textbf{Nine modules are genuinely exercised and all
nine are consistent with conformance.} Two, request-uri-method-post and
request-uri-fetched-twice, test fetching the request by reference; the suite's own
descriptions say they do not apply to \texttt{request\_method=url\_query}, so under our variant
they ran as plain happy flows. They are recorded as \emph{not exercised}, not as passes.
Exercising them needs a signed by-reference request with \texttt{x509\_san\_dns}, present only on
the \texttt{release-1.0.x} head.

The 1.0 baseline now covers all 11 modules, each with its reason and evidence; re-gating this
run against it passes with no unknown modules and no mismatches.

\subsubsection*{Status at end of Day 6}
The verifier plan is complete on the target version: 9 of 11 modules exercised, all
consistent with conformance; 2 not applicable to our variant. Earlier in the day, before the
full run: one module of twelve passed. Two new findings (F-10, F-11); F-05
resolved as our artefact; F-01, F-02, F-03, F-04 confirmed fixed on 1.0 by run; F-07
persists on 1.0; F-09 does not apply on 1.0. Next: Inji Certify. Question for mentors, sharpened: is \texttt{1.0.0-alpha.1} the target,
or the \texttt{release-1.0.x} head built from source?

\newpage
%=========================================================
\section{Findings Catalogue}

Severity reflects impact on certifiability and on security, not on our convenience.

\textbf{Version scope (added 2 October).} Every finding below was established against
\textbf{Inji Verify 0.18.2}. On 2 October we found an Inji 1.0 line (\texttt{release-1.0.x};
runnable as \texttt{1.0.0-alpha.1}) in which F-01, F-02 and F-03 appear fixed from source and
F-09 is partly addressed. The findings stand for 0.18.2; none is to be reported upstream
until it has been re-checked against the 1.0 line. See the Day 6 entry. F-10 and F-11 were
found on the 1.0 line itself.

%---------------------------------------------------------
\finding{F-01 --- Nonce has insufficient entropy}{\sevhigh{} \quad \confirmed}

\textbf{Component:} Inji Verify 0.18.2 \quad \textbf{Spec:} OID4VP 1.0 Final, §5.2

\textbf{Observation.} The \texttt{nonce} in Verify's OID4VP authorisation request is a
base64-encoded millisecond Unix timestamp rather than a cryptographically random value.

\textbf{Evidence --- our analysis.} Three samples decoded, across both request-construction
paths:
\begin{longtable}{p{3.5cm} p{4.5cm} p{3cm} p{2.5cm}}
\hline
\textbf{clientIdScheme} & \textbf{nonce (base64)} & \textbf{decoded} & \textbf{$\Delta$ to issuedAt}\\
\hline
\endhead
\texttt{did} & \texttt{MTc4OTg4MzIxMzQ5MA==} & 1789883213490 & ---\\
\texttt{pre\_registered} & \texttt{MTc4OTg4NDUzNjY3OQ==} & 1789884536679 & 120\,ms\\
\texttt{pre\_registered} & \texttt{MTc4OTg4NDY1NjQzNg==} & 1789884656436 & 158\,ms\\
\hline
\end{longtable}

\textbf{Evidence --- conformance suite, independently.}
\begin{quote}\small
\texttt{VP1FinalEnsureMinimumNonceEntropy} --- WARNING\\
``Nonce does not appear to have sufficient entropy (calculated Shannon entropy is below the
threshold). Section 5.2 of OID4VP 1.0 Final requires the nonce to be a `fresh,
cryptographically random number with sufficient entropy'; OpenID4VP PR \#722 recommends at
least 128 bits.''\\
\texttt{actual: 73.68} \quad \texttt{expected: 96.0}
\end{quote}

\textbf{Root cause --- located 23 September, and it corrects our Day 1 attribution.}
The weak nonce is generated by the \emph{SDK}, not by the service. In
\texttt{inji-verify-sdk/src/utils/api.ts} lines 10--12:

\begin{quote}\small
\texttt{const generateNonce = (): string => \{}\\
\texttt{\ \ return btoa(Date.now().toString());}\\
\texttt{\};}
\end{quote}

\noindent used at \texttt{api.ts:89} and \texttt{api.ts:159}. \texttt{verify-ui} depends on
\texttt{@injistack/react-inji-verify-sdk} 0.18.1, so this is the nonce that reaches the
authorisation request. \texttt{btoa} is standard base64, which pads with \texttt{=} ---
the same line therefore causes F-02.

\textbf{The service already implements this correctly.}
\texttt{verify-service/.../utils/SecurityUtils.java} generates a nonce from
\texttt{SecureRandom} over 16 bytes. But at
\texttt{VerifiablePresentationRequestServiceImpl.java:86} it is only a fallback:
\texttt{vpRequestCreate.getNonce() != null ? vpRequestCreate.getNonce() :
SecurityUtils.generateNonce()}. Because the SDK always supplies a nonce, the sound
generator never executes in the normal flow.

This reframes the finding materially: it is \emph{not} a missing capability but a correct
implementation overridden by its caller, which makes it a small, well-scoped fix rather
than a design change --- and a much more actionable report for a maintainer. Our Day 1
write-up attributed it to Inji Verify generally, which was imprecise. Recorded as a
reminder that the ``read the source, don't guess'' rule applies to root causes as much as
to behaviour.

\textbf{Confirmed experimentally --- 29 September.} Two automated runs differing only in
the nonce source (\texttt{hcuWEvxGbvRgt7e}, SDK-shaped; \texttt{Q9MaDZyfpJhjiDp}, service-generated):
this check and F-02's move from WARNING/FAILURE to SUCCESS, and nothing else changes. See
the Day 4 entry.

\textbf{Impact.} The nonce exists to bind a wallet's presentation to a specific request and
prevent replay. A value derived from the current clock is predictable to within milliseconds,
substantially weakening that guarantee. In a national identity deployment this is a
meaningful security property, not a cosmetic issue.

\textbf{ON HOLD --- 2 October.} The Inji 1.0 line already generates the nonce with
\texttt{crypto.getRandomValues}. The issue below must not be filed as drafted; whether it is
worth raising for the 0.18.x line is a question for the mentors.

\textbf{Proposed action (superseded).} File upstream against \texttt{mosip/inji-verify} as a single issue
covering F-01 and F-02 together, since they share one root cause and one fix. Recommend
\texttt{crypto.getRandomValues} over 32 bytes, base64url-encoded without padding. Draft
prepared at \texttt{docs/upstream/ISSUE-01-nonce-entropy.md}; not yet filed.

%---------------------------------------------------------
\finding{F-02 --- Nonce contains non-URL-safe characters}{\sevmed{} \quad \confirmed}

\textbf{Component:} Inji Verify 0.18.2 \quad \textbf{Spec:} OID4VP 1.0 Final, §5.2

\textbf{Evidence.}
\begin{quote}\small
\texttt{CheckForInvalidCharsInNonce} --- FAILURE\\
``Non URL safe characters found in nonce. This may introduce interoperability issues.''\\
\texttt{invalid\_chars: ["="]}
\end{quote}

\textbf{Observation.} The nonce uses standard base64 with \texttt{=} padding rather than
base64url. The \texttt{=} character is not URL-safe and requires percent-encoding when the
nonce travels in a query string, which not all wallets handle consistently.

\textbf{Root cause.} Identical to F-01 --- \texttt{btoa} in
\texttt{inji-verify-sdk/src/utils/api.ts:11} emits standard base64 with \texttt{=} padding
rather than base64url.

\textbf{Impact.} Interoperability risk with third-party wallets. Same root cause as F-01 and
fixed by the same change; the two are filed upstream as one issue.

%---------------------------------------------------------
\finding{F-03 --- No DCQL support; test plan cannot proceed}{\sevhigh{} \quad \confirmed}

\textbf{Component:} Inji Verify 0.18.2 \quad \textbf{Spec:} OID4VP 1.0 Final, §6

\textbf{Evidence --- conformance suite.}
\begin{quote}\small
\texttt{ExtractDCQLQueryFromAuthorizationRequest} --- FAILURE\\
``dcql\_query not found in authorization request parameters''
\end{quote}

\textbf{Evidence --- source.} A case-insensitive search for \texttt{dcql} across
\texttt{verify-service/src/main} returns zero matches. Credential selection is implemented
entirely via \texttt{presentation\_definition} with \texttt{input\_descriptors} (DIF
Presentation Exchange), set in \texttt{VerifiablePresentationRequestServiceImpl} as the
\texttt{presentation\_definition\_uri} claim.

\textbf{Impact.} This terminates the test. Because every module in the 1.0 Final verifier
plan requires a DCQL query, \textbf{Inji Verify 0.18.2 cannot pass any test in that plan.}
This is the single most consequential finding of Day 1 and is not fixable by configuration.

\textbf{Open question for mentors.} See Q1.

%---------------------------------------------------------
\finding{F-04 --- Uses superseded \texttt{presentation\_definition}}{\sevmed{} \quad \confirmed}

\textbf{Component:} Inji Verify 0.18.2 \quad \textbf{Spec:} OID4VP 1.0 Final

\textbf{Evidence.}
\begin{quote}\small
\texttt{CheckNoPresentationDefinitionInVpAuthorizationRequest} --- WARNING\\
``Authorization request contains legacy `presentation\_definition'. VP1 Final uses
`dcql\_query' for credential selection, so sending both is unnecessary and may indicate
draft and final request formats are being mixed.''
\end{quote}

\textbf{Relationship to F-03.} Two faces of the same gap: Verify implements the earlier
draft's query model. Recorded separately because a maintainer may address them
independently (e.g.\ support both during a transition period).

%---------------------------------------------------------
\finding{F-05 --- \texttt{client\_metadata} absent in inline mode}{\sevlow{} \quad \textbf{RESOLVED --- NOT A DEFECT}}

\textbf{Resolved 2 October: our artefact.} The 0.18.x SDK adds \texttt{client\_metadata}
(\texttt{client\_name} and \texttt{vp\_formats}) to every inline request when it builds the
request URI (\texttt{OpenID4VPVerification.tsx}, lines 124--130). Our
\texttt{build-oid4vp-uri.py}, and the runner copied from it, did not. The ``client\_metadata
not present'' warning in the Day 1 run and in run \texttt{hcuWEvxGbvRgt7e} was produced by our
tooling, and both runs share the artefact --- which is why they matched each other. The runner
now reproduces the SDK. One caveat stands: the 0.18 SDK uses the pre-Final key
\texttt{vp\_formats}, while 1.0 Final requires \texttt{vp\_formats\_supported}, so on 0.18.2 that
check would likely still fail, for a different reason. On 1.0 the SDK uses the Final key.

\textbf{Original entry follows.}

\textbf{Component:} Inji Verify 0.18.2 \quad \textbf{Spec:} OID4VP 1.0 Final, §5.1; Annex B.2.2, B.3.4

\textbf{Evidence.}
\begin{quote}\small
\texttt{VP1FinalValidateVpFormatsSupportedInClientMetadata} --- FAILURE\\
``client\_metadata is missing or not a JSON object''\\[0.4em]
\texttt{VP1FinalCheckForUnexpectedParametersInVpClientMetadata} --- WARNING\\
``client\_metadata not present in request but it is mandatory to include vp\_formats in it''
\end{quote}

\textbf{Why this is not yet confirmed.} In the \texttt{did} (by-reference) path, Verify's
signed request object \emph{did} contain \texttt{client\_metadata} with a full
\texttt{vp\_formats} block advertising \texttt{ldp\_vp} and \texttt{vc+sd-jwt}. In the
\texttt{pre\_registered} (inline) path, the \texttt{/vp-session-request} API response
contained no \texttt{client\_metadata} field, so the URI we constructed could not include
one. It is not yet established whether Verify omits it in inline mode, or whether the UI
adds it when building the QR code and our reconstruction simply dropped it.

\textbf{Action required before reporting.} Inspect the QR payload the Verify UI actually
produces in \texttt{pre\_registered} mode and compare against the API response.

%---------------------------------------------------------
\finding{F-06 --- No \texttt{direct\_post.jwt}; HAIP certification impossible}{\sevhigh{} \quad \confirmed}

\textbf{Component:} Inji Verify 0.18.2 \quad \textbf{Spec:} OID4VP 1.0 Final / HAIP

\textbf{Evidence --- source.}
\begin{quote}\small
\texttt{verify-service/src/main/java/io/inji/verify/shared/Constants.java}\\
\texttt{public static final String RESPONSE\_MODE = "direct\_post";}
\end{quote}
This is a \texttt{static final} constant, not a configuration property or environment
variable. No code path produces the encrypted \texttt{direct\_post.jwt} variant.

\textbf{Evidence --- conformance suite.} The HAIP verifier plan
(\emph{OpenID for Verifiable Presentations 1.0 Final/HAIP: Test a verifier}) offers only
\texttt{direct\_post.jwt} in its Response Mode variant.

\textbf{Impact.} Inji Verify cannot be submitted for HAIP self-certification in its current
state. HAIP is the profile relevant to high-assurance national identity deployments, which
is MOSIP's own domain.

\textbf{Open question for mentors.} See Q2.

%---------------------------------------------------------
\finding{F-07 --- DID document verification method identifier inconsistency}{\sevmed{} \quad \confirmed}

\textbf{Component:} Inji Verify 0.18.2 \quad \textbf{Spec:} W3C DID Core

\textbf{Evidence.} Retrieved from
\texttt{https://<host>/v1/verify/did.json}:
\begin{quote}\small
\texttt{"id": "did:web:<host>:v1:verify"} \quad (document subject)\\
\texttt{"verificationMethod[0].id": "did:web:<host>:verify\#key-0"} \quad (no \texttt{:v1:})\\
\texttt{"assertionMethod[0]": "did:web:<host>:verify\#key-0"} \quad (no \texttt{:v1:})
\end{quote}

\textbf{Origin.} The default Docker Compose environment variables differ:
\texttt{INJI\_DID\_VERIFY\_URI} includes the \texttt{v1} path segment while
\texttt{INJI\_DID\_VERIFY\_PUBLIC\_KEY\_URI} does not.

\textbf{Impact.} Under \texttt{did:web} resolution rules, the key's identifier resolves to
\texttt{https://<host>/verify/did.json}, a path that does not exist on the deployment. A
strict relying party resolving the verification method by its own identifier would fail.

%---------------------------------------------------------
\finding{F-08 --- No mdoc / ISO 18013-5 support}{\sevlow{} \quad \confirmed}

\textbf{Component:} Inji Verify 0.18.2

\textbf{Evidence.} Searches for \texttt{mso\_mdoc}, \texttt{iso\_mdl} and \texttt{mdoc}
across \texttt{verify-service/src/main} return zero matches.

\textbf{Impact.} The \texttt{iso\_mdl} credential format variant of the verifier plan cannot
be exercised. Recorded for completeness; likely a deliberate scope decision rather than a
defect.

%---------------------------------------------------------
\finding{F-09 --- Client identifier prefix has no overlap with 1.0 Final}{\sevmed{} \quad \confirmed}

\textbf{Component:} Inji Verify 0.18.2 \quad \textbf{Spec:} OID4VP 1.0 Final, §5.9.2

\textbf{Evidence.}
\begin{quote}\small
\texttt{EnsureClientIdMatchesResponseUri} --- FAILURE\\
``For redirect\_uri client\_id prefix, the client\_id (without prefix) must equal
response\_uri''\\
\texttt{client\_id: "inji-verify-ui"} \quad
\texttt{response\_uri: "https://<host>/v1/verify/vp-submission/direct-post"}
\end{quote}

\textbf{Observation.} Inji Verify supports two client identifier schemes, \texttt{did} and
\texttt{pre\_registered}. The 1.0 Final verifier plan offers three prefixes:
\texttt{redirect\_uri}, \texttt{x509\_san\_dns}, \texttt{x509\_hash}. Verify has no X.509
material configured, so the two X.509 options are unreachable, and \texttt{redirect\_uri}
requires a client identifier Verify does not produce. \textbf{No variant combination in the
plan matches Verify's client identification model.}

\textbf{Caveat.} This failure is partly a consequence of the variant we selected. The
reportable observation is the absence of overlap, not that any single value is ``wrong''.

\textbf{Open question for mentors.} See Q3.

%---------------------------------------------------------
\finding{F-10 --- Docker Compose database script out of date; every VP submission fails with HTTP 500}{\sevhigh{} \quad \confirmed}

\textbf{Component:} Inji Verify 1.0.0-alpha.1, and \texttt{release-1.0.x} head (6a6ab3b, 1 October)
\quad \textbf{Kind:} packaging --- the service is correct

\textbf{Observation.} Brought up the documented way, with Docker Compose, Inji Verify 1.0
answers every presentation submission (\texttt{POST /v2/vp-submission/direct-post}) with HTTP 500.

\textbf{Evidence --- service log} (run \texttt{gIA512EeJWehcHt}):
\begin{quote}\small
\texttt{PSQLException: ERROR: column v1\_0.response\_code does not exist}\\
\texttt{at ...VerifiablePresentationRequestServiceImpl.getCurrentRequestStatus(...java:126)}
\end{quote}

\textbf{Root cause.} The Compose stack initialises its database from
\texttt{docker-compose/db-init/init.sql}, which creates \texttt{verify.vp\_submission} with five
columns. The canonical schema, \texttt{db\_scripts/inji\_verify/ddl/verify-vp\_submission.sql},
has eight: it adds \texttt{response\_code}, \texttt{response\_code\_expiry\_at} and
\texttt{response\_code\_used}, makes \texttt{vp\_token} and \texttt{presentation\_submission}
nullable, and adds a primary key. The 1.0 service is written against the canonical schema.

\textbf{Ruled out first:} a stale database from the earlier 0.18.2 stack. The data volume was
created at 06:34 on 2 October; the Postgres log shows \texttt{init.sql} running and creating
four tables; the live table matches \texttt{init.sql} column for column.

\textbf{Confirmed fix.} With the canonical \texttt{vp\_submission} definition substituted into
\texttt{init.sql}, the 500 disappears and Verify processes submissions normally (runs
\texttt{tpKVz6upccxpO5U} onward).

\textbf{Also inferred, not observed:} \texttt{presentation\_submission NOT NULL} would likely
break storing a DCQL response, which carries no presentation submission. Not separately tested.

\textbf{Impact.} Anyone evaluating Inji Verify 1.0 through its quick start sees verification
fail outright, for a reason that is not in the product.

\textbf{Proposed action.} Report upstream with the fix: regenerate \texttt{init.sql} from
\texttt{db\_scripts}. Present on the branch head, so this is the one to file. Draft in
\texttt{docs/upstream/ISSUE-02-compose-db-init.md}.

%---------------------------------------------------------
\finding{F-11 --- Shipped DCQL sample uses the superseded format identifier}{\sevlow{} \quad \confirmed}

\textbf{Component:} Inji Verify 1.0.0-alpha.1 \texttt{docker-compose/config/config.json}
\quad \textbf{Kind:} sample configuration

\textbf{Evidence.} \texttt{ValidateDCQLQuery}: \emph{``does not have a value in the enumeration
[dc+sd-jwt, mso\_mdoc]''} at \texttt{\$.credentials[0].format}; the shipped \emph{Mock Identity}
entry requests \texttt{vc+sd-jwt}. OID4VP 1.0 Final, §6.

\textbf{Configuration, not implementation.} Run \texttt{ZQqnbYfw46wpmyI}, requesting
\texttt{dc+sd-jwt}, cleared both DCQL failures and Verify accepted the presentation.

\textbf{Not to be reported.} The \texttt{release-1.0.x} head already adds a \texttt{dc+sd-jwt}
entry to the sample configuration; upstream is addressing it.

\newpage
%=========================================================
\section{Open Questions for Mentors}

These are the questions we would like to put to MOSIP mentors. Each is phrased to be
answerable without our setup in front of them.

\textbf{Note on channel.} The kickoff seminar and mentor AMA scheduled for the start of the
event did not take place, so the community forum is currently the only route to an answer,
and no answer is guaranteed. Where a question shapes the build --- Q1 above all --- our
approach is to make the affected choice a configuration value and to state the assumption
plainly in the submission, rather than to block on a reply.

\begin{description}

\item[\textbf{Q1 --- Target specification version.}]\mbox{}\\
Inji Verify 0.18.2 implements OID4VP credential selection via \texttt{presentation\_definition}
(Presentation Exchange), whereas OID4VP 1.0 Final requires \texttt{dcql\_query}. Consequently
no test in the 1.0 Final verifier plan can complete. Which specification version should our
harness target as the benchmark: 1.0 Final (accepting that the current baseline is a full
set of failures), or the ID2 draft that the shipped implementation appears to target? Is
DCQL support on the Inji roadmap, and if so for which release?

\item[\textbf{Q2 --- HAIP certification intent.}]\mbox{}\\
The problem statement asks us to document ``how to submit for self-certification''. The HAIP
verifier plan mandates \texttt{direct\_post.jwt}, but \texttt{RESPONSE\_MODE} is a
\texttt{static final} constant set to \texttt{direct\_post} in \texttt{Constants.java}. Is
encrypted response mode planned? Should our harness include the HAIP plan as a
known-failing suite to track progress toward certification, or exclude it until the
capability exists?

\item[\textbf{Q3 --- Client identifier scheme.}]\mbox{}\\
Verify supports \texttt{did} and \texttt{pre\_registered} client identifier schemes; the 1.0
Final plan offers \texttt{redirect\_uri}, \texttt{x509\_san\_dns} and \texttt{x509\_hash}.
Is there a supported configuration that allows an Inji Verify deployment to present a
client identifier the 1.0 Final conformance plan can validate, or is this a known gap?

\item[\textbf{Q4 --- Benchmark semantics for \texttt{api-testrig}.}]\mbox{}\\
The problem statement asks for results ``gated against a conformance benchmark''. Given that
the honest current baseline is ``all modules fail at DCQL extraction'', should the benchmark
be (a) a frozen snapshot of today's pass/fail set, so the gate catches regressions only, or
(b) an aspirational target that the build fails against until conformance improves? Our
current design assumes (a), with an expected-failures file under version control.

\item[\textbf{Q5 --- Sample configuration shipped with Inji Verify.}]\mbox{}\\
The default \texttt{config.json} ships five credential types, four of which use
\texttt{ldp\_vc} with \texttt{clientIdScheme: did}. That combination cannot be exercised by
the 1.0 Final plan at all. Is the shipped sample configuration intended to be
certification-representative, or purely a demonstration fixture? This affects whether we
report the configuration itself as a finding.

\item[\textbf{Q6 --- Interactive tests and automation.}]\mbox{}\\
Every module in the verifier plan is interactive: it waits for a human to trigger a flow in
the verifier and paste the resulting authorisation request, and the happy-flow module also
requires a screenshot upload before it reaches REVIEW status. Our harness can drive request
generation via Verify's own API, but the screenshot-upload requirement appears to need
\texttt{POST /api/log/\{id\}/images}. Is a fully unattended run of the verifier plan
considered achievable, or is a documented semi-automated flow the expected deliverable?

\textbf{Substantially answered by our own investigation on Day 3}, and retained here only
for confirmation. The authorisation request is delivered by an HTTP GET to the test's exposed
\texttt{authorization\_endpoint}, and the screenshot placeholder is filled by the suite's own
\texttt{browser} automation against an HTML stand-in page it serves itself. Both are
source-derived and await confirmation against a live run.

\item[\textbf{Q7 --- Scope of Inji Certify work.}]\mbox{}\\
The \texttt{docker-compose-injistack} setup requires a PKCS12 keystore obtained through
Mimoto onboarding, a publicly resolvable DID endpoint, and an external authorisation server.
For a two-person team on a four-week timeline, is there a lighter-weight Certify
configuration suitable for conformance testing, or is the Collab sandbox the intended target
for the issuer plan?

\item[\textbf{Q8 --- Presentation definition id (raised 29 September).}]\mbox{}\\
\texttt{verify-ui} sends every presentation definition with the same hard-coded id
(\texttt{vpVerificationState.ts:59}), while the definition's contents vary with the credentials
selected. Is anything on the Inji side --- or in Inji Wallet --- expected to rely on that id
distinguishing definitions? We have not treated this as a defect, because we have found no
consequence of it.

\end{description}

\newpage
%=========================================================
\section{Reproduction}

To reproduce the Day 1 baseline run:

\begin{enumerate}\tightlist
  \item Follow \texttt{docs/SETUP.md} to install prerequisites and clone the repositories.
  \item Follow \texttt{docs/RUNBOOK.md} to start the conformance suite, Inji Verify, and ngrok.
  \item In \texttt{inji-verify/docker-compose/config/config.json}, ensure the
        \emph{Mock Identity (SD JWT)} entry has \texttt{"clientIdScheme":"pre\_registered"}.
  \item In the conformance suite, create the plan
        \emph{OpenID for Verifiable Presentations 1.0 Final: Test a verifier --- alpha tests}
        with the variant recorded in Section 2.
  \item Run module 01, \texttt{oid4vp-1final-verifier-happy-flow}.
  \item In the Verify UI, request the \emph{Mock Identity (SD JWT)} credential; capture the
        \texttt{/v1/verify/vp-session-request} response.
  \item Run \texttt{python3 scripts/build-oid4vp-uri.py} on that response and paste the
        resulting URI into the suite. The request expires 300 seconds after issue.
  \item Export the signed log and store it under \texttt{logs/<date>/}.
\end{enumerate}

\vspace{1em}
\noindent\rule{\textwidth}{0.4pt}\\[0.3em]
\small\textit{Since Day 4 the whole reproduction above is one command:}
\texttt{./run-conformance.sh --component verify}.
\textit{Findings from modules 02--11 and the ID2 verifier plan to be appended as they are executed.}

\end{document}
